\documentclass[11pt,a4paper]{article}

\usepackage[utf8]{inputenc}
\usepackage[T1]{fontenc}
\usepackage[french]{babel}
\usepackage[margin=2.4cm,top=2.6cm,bottom=2.4cm]{geometry}
\usepackage{lmodern}
\usepackage{microtype}
\usepackage{booktabs}
\usepackage{array}
\usepackage{enumitem}
\usepackage{listings}
\usepackage{xcolor}
\usepackage{titlesec}
\usepackage{fancyhdr}
\usepackage{tcolorbox}
\usepackage{tikz}
\usetikzlibrary{positioning,arrows,calc,fit,backgrounds}
\usepackage[colorlinks=true,linkcolor=accent,urlcolor=accent,citecolor=accent]{hyperref}

% ---------------------------------------------------------------- palette ----
\definecolor{accent}{HTML}{4F46E5}
\definecolor{ink}{HTML}{18181B}
\definecolor{muted}{HTML}{71717A}
\definecolor{hairline}{HTML}{E4E4E7}
\definecolor{surface}{HTML}{FAFAFA}
\definecolor{warnbg}{HTML}{FFFBEB}
\definecolor{warnline}{HTML}{D97706}

\color{ink}

% ------------------------------------------------------------- structure ----
\titleformat{\section}{\large\bfseries\color{accent}}{\thesection}{0.6em}{}
\titleformat{\subsection}{\normalsize\bfseries\color{ink}}{\thesubsection}{0.6em}{}
\titlespacing*{\section}{0pt}{1.4em}{0.5em}
\titlespacing*{\subsection}{0pt}{1.0em}{0.35em}

\pagestyle{fancy}
\fancyhf{}
\renewcommand{\headrulewidth}{0.4pt}
\renewcommand{\headrule}{\color{hairline}\hrule height\headrulewidth}
\fancyhead[L]{\footnotesize\color{muted}Transcription vocale temps réel — contrat d'intégration}
\fancyhead[R]{\footnotesize\color{muted}\thepage}

% ----------------------------------------------------------------- code ----
\lstdefinestyle{docie}{
  basicstyle=\ttfamily\small,
  backgroundcolor=\color{surface},
  frame=single,
  rulecolor=\color{hairline},
  framesep=6pt,
  xleftmargin=2pt,
  xrightmargin=2pt,
  breaklines=true,
  columns=fullflexible,
  keepspaces=true,
  showstringspaces=false,
  stringstyle=\color{accent},
  commentstyle=\color{muted}\itshape,
  aboveskip=0.9em,
  belowskip=0.9em,
}
\lstset{style=docie}

% ------------------------------------------------------------ tikz style ----
\tikzset{
  lane/.style={draw=hairline, fill=white, rounded corners=3pt, line width=0.7pt,
               minimum width=3.3cm, minimum height=0.95cm, align=center,
               font=\small\bfseries},
  life/.style={draw=hairline, dashed, line width=0.7pt},
  msg/.style={->, >=latex, line width=0.8pt, draw=accent, font=\scriptsize},
  msgback/.style={->, >=latex, line width=0.8pt, draw=muted, font=\scriptsize, dashed},
  audio/.style={->, >=latex, line width=1.6pt, draw=accent, font=\scriptsize},
  stepnum/.style={circle, fill=accent, text=white, font=\scriptsize\bfseries,
               inner sep=1.4pt, minimum size=3.6mm},
  box/.style={draw=hairline, fill=white, rounded corners=3pt, line width=0.7pt,
              minimum height=1.05cm, align=center, font=\small},
  note/.style={font=\scriptsize\itshape, text=muted, align=center},
}

% --------------------------------------------------------------- boxes ----
\newtcolorbox{statusbox}{
  colback=warnbg, colframe=warnline, boxrule=0.8pt, arc=3pt,
  left=10pt, right=10pt, top=8pt, bottom=8pt,
}
\newtcolorbox{dontbox}[1]{
  colback=white, colframe=warnline, boxrule=1.2pt, arc=3pt,
  left=10pt, right=10pt, top=8pt, bottom=8pt,
  title={\bfseries #1}, coltitle=white, colbacktitle=warnline,
  fonttitle=\normalsize,
}
\newtcolorbox{askbox}{
  colback=white, colframe=accent, boxrule=1pt, arc=3pt,
  left=10pt, right=10pt, top=8pt, bottom=8pt,
}

\setlist[itemize]{leftmargin=1.3em, itemsep=0.32em, topsep=0.35em}
\setlist[enumerate]{leftmargin=1.5em, itemsep=0.45em, topsep=0.35em}

\begin{document}
\thispagestyle{empty}

% ----------------------------------------------------------- title block ----
\begin{flushleft}
{\color{accent}\rule{\textwidth}{2.2pt}}\\[1.1em]
{\Huge\bfseries Transcription vocale temps réel}\\[0.45em]
{\Large\color{muted} Contrat d'intégration}\\[1.1em]
{\color{hairline}\rule{\textwidth}{0.6pt}}\\[0.7em]
{\small\color{muted}
Plateforme d'inférence DocIE \quad$\longleftrightarrow$\quad ERP Cynov \hfill \today}
\end{flushleft}

\vspace{1.3em}

\begin{statusbox}
\textbf{Statut : proposition soumise à relecture. Non encore implémenté.}\par\smallskip
Ce document existe pour que les deux équipes valident le contrat \emph{avant} que l'une
ou l'autre ne développe. Les chemins d'accès, les noms de champs et les types de
messages peuvent encore évoluer selon vos retours. Merci de ne pas démarrer les
développements avant validation conjointe de la section~\ref{sec:erp}.
\end{statusbox}

\vspace{0.6em}

\section{Ce que nous construisons}

L'ERP capte le micro dans le navigateur et affiche une transcription en direct pendant
que l'utilisateur parle. L'audio transite vers la plateforme d'inférence DocIE via une
WebSocket ; les transcriptions partielles reviennent en une seconde environ, et une
transcription finale stable est émise à chaque pause du locuteur.

Le modèle gère le français (fr-FR et fr-CA) et diffuse nativement en continu : la
transcription est donc incrémentale, et non du type \emph{envoyer un fichier puis
attendre}. La transcription de fichiers est explicitement hors périmètre.

\vspace{0.8em}
\begin{center}
\begin{tikzpicture}[node distance=0pt]

  \node[lane] (be)  at (0,0)     {Backend ERP};
  \node[lane] (br)  at (5.6,0)   {Navigateur ERP};
  \node[lane] (dc)  at (11.2,0)  {Plateforme DocIE};

  \draw[life] (be.south)  -- ++(0,-6.25);
  \draw[life] (br.south)  -- ++(0,-6.25);
  \draw[life] (dc.south)  -- ++(0,-6.25);

  \draw[msg] ([yshift=-1.05cm]be.south) -- node[above,midway,align=center]
       {\texttt{POST /v1/realtime/sessions}\\ en-tête \texttt{X-API-Key}}
       ([yshift=-1.05cm]dc.south);
  \node[stepnum] at ([xshift=-3mm,yshift=-1.05cm]be.south) {1};

  \draw[msgback] ([yshift=-2.15cm]dc.south) -- node[above,midway]
       {\texttt{\{ ticket, expires\_in: 60 \}}}
       ([yshift=-2.15cm]be.south);
  \node[stepnum] at ([xshift=3mm,yshift=-2.15cm]dc.south) {2};

  \draw[msg] ([yshift=-3.15cm]be.south) -- node[above,midway]
       {ticket transmis à sa propre page}
       ([yshift=-3.15cm]br.south);
  \node[stepnum] at ([xshift=-3mm,yshift=-3.15cm]be.south) {3};

  \draw[msg] ([yshift=-4.15cm]br.south) -- node[above,midway,align=center]
       {\texttt{WSS /v1/realtime}\\ ticket en sous-protocole}
       ([yshift=-4.15cm]dc.south);
  \node[stepnum] at ([xshift=-3mm,yshift=-4.15cm]br.south) {4};

  \draw[audio] ([yshift=-5.35cm]br.south) -- node[above,midway]
       {PCM16 16\,kHz mono, trames binaires}
       ([yshift=-5.35cm]dc.south);
  \node[stepnum] at ([xshift=-3mm,yshift=-5.35cm]br.south) {5};

  \draw[msgback] ([yshift=-6.05cm]dc.south) -- node[below,midway]
       {\texttt{partial} \ldots\ \texttt{partial} \ldots\ \texttt{final}}
       ([yshift=-6.05cm]br.south);
  \node[stepnum] at ([xshift=3mm,yshift=-6.05cm]dc.south) {6};

\end{tikzpicture}

\vspace{0.3em}
{\scriptsize\color{muted} Les étapes 1 à 3 sont de serveur à serveur et n'ont lieu qu'une
fois par session de dictée. Seule la flèche épaisse transporte de l'audio.}
\end{center}

\section{Le point à ne surtout pas manquer}

\begin{dontbox}{Ne jamais transmettre une clé API DocIE au navigateur}
Les clés API de notre plateforme ont une durée de vie longue et une portée totale : une
seule clé donne accès à tous les modèles de l'API \texttt{/v1}, et toute la consommation
lui est imputée et auditée. Une clé intégrée au JavaScript de l'ERP est lisible par
chaque utilisateur de l'ERP et par tout ce qui s'exécute dans cette page.\par\smallskip
Traitez-la exactement comme un mot de passe de base de données.
\end{dontbox}

Ce n'est pas une restriction que nous pouvons lever par configuration, et c'est la raison
pour laquelle le flux comporte deux étapes au lieu d'une. Un navigateur n'est pas un
endroit où peut résider une authentification durable : il reçoit donc un ticket éphémère.
La clé API ne quitte jamais votre serveur.

\section{Ce que l'équipe ERP implémente}
\label{sec:erp}

Deux développements, un de chaque côté de votre pile.

\subsection{Backend : émettre un ticket (serveur à serveur)}

Votre backend détient déjà la clé API DocIE. Lorsqu'un utilisateur ouvre l'interface de
dictée, votre backend demande un ticket et le transmet à la page de cet utilisateur.

\begin{lstlisting}
curl -X POST https://api.inference.cynov.com/v1/realtime/sessions \
  -H "X-API-Key: $DOCIE_API_KEY" \
  -H "Content-Type: application/json" \
  -d '{"language": "fr"}'
\end{lstlisting}

\begin{lstlisting}
{
  "ticket":     "docie_rt_<opaque>",
  "expires_in": 60,
  "ws_url":     "wss://api.inference.cynov.com/v1/realtime"
}
\end{lstlisting}

\noindent Propriétés du ticket, qui déterminent son usage :

\begin{center}
\renewcommand{\arraystretch}{1.35}
\begin{tabular}{@{}>{\bfseries}p{3.4cm}p{10.9cm}@{}}
\toprule
Éphémère      & Environ 60 secondes : de quoi ouvrir la socket, pas de quoi mériter
                d'être stocké. À émettre au moment où l'utilisateur lance la dictée. \\
Usage unique  & Consommé par la première socket qui le présente ; une seconde connexion
                avec le même ticket est refusée. \\
Portée réduite & Ouvre une socket de transcription et n'autorise rien d'autre. Ce n'est
                pas un substitut à la clé API. \\
Imputé        & Toute la consommation est enregistrée sur la clé API émettrice : votre
                suivi reste donc par clé, quel que soit le nombre d'utilisateurs. \\
\bottomrule
\end{tabular}
\end{center}

\noindent Ne mettez pas les tickets en cache et n'en partagez pas un entre plusieurs
utilisateurs. L'émission est peu coûteuse.

\subsection{Navigateur : ouvrir la socket et diffuser l'audio}

\begin{lstlisting}
const ws = new WebSocket(ws_url, [ticket]);   // ticket en sous-protocole
ws.binaryType = "arraybuffer";
\end{lstlisting}

Le ticket voyage comme sous-protocole WebSocket (\texttt{Sec-WebSocket-Protocol}), le seul
en-tête qu'un navigateur peut positionner lors d'une poignée de main WebSocket. Nous le
gardons délibérément hors de la chaîne de requête, car les URL finissent dans les journaux
des proxys, les journaux d'accès et l'historique du navigateur. Si votre pile ne peut pas
définir de sous-protocole, dites-le-nous et nous activerons un repli par paramètre
d'URL — mais le sous-protocole reste préférable.

\vspace{0.9em}
\begin{center}
\begin{tikzpicture}[node distance=5.5mm]

  \node[box, minimum width=2.15cm] (mic)  {Microphone\\[-1pt]{\scriptsize 48\,kHz}};
  \node[box, minimum width=2.45cm, right=of mic] (work) {\texttt{AudioWorklet}\\[-1pt]{\scriptsize sous-échantillonnage}};
  \node[box, minimum width=2.15cm, right=of work] (pcm)  {PCM16\\[-1pt]{\scriptsize 16\,kHz mono}};
  \node[box, minimum width=2.15cm, right=of pcm, fill=surface] (ws)   {WebSocket\\[-1pt]{\scriptsize trames binaires}};
  \node[box, minimum width=2.45cm, right=of ws] (asr)  {ASR en continu\\[-1pt]{\scriptsize blocs 320--560\,ms}};
  \node[box, minimum width=2.15cm, right=of asr] (out)  {Transcription\\[-1pt]{\scriptsize partielle / finale}};

  \foreach \a/\b in {mic/work, work/pcm, pcm/ws, ws/asr, asr/out}
    \draw[->, >=latex, line width=0.8pt, draw=accent] (\a) -- (\b);

  \node[note, below=3.5mm of pcm, xshift=-6mm] {dans la page ERP};
  \node[note, below=3.5mm of asr, xshift=4mm] {sur la plateforme d'inférence};

  \begin{scope}[on background layer]
    \node[draw=hairline, dashed, rounded corners=4pt, inner sep=3.6mm,
          fit=(mic)(pcm)] {};
    \node[draw=hairline, dashed, rounded corners=4pt, inner sep=3.6mm,
          fit=(asr)(out)] {};
  \end{scope}

\end{tikzpicture}
\end{center}

\subsubsection*{Format audio}

\begin{center}
\renewcommand{\arraystretch}{1.25}
\begin{tabular}{@{}>{\bfseries}ll@{}}
\toprule
Encodage          & PCM signé 16 bits, petit-boutiste \\
Fréquence         & 16\,000 Hz \\
Canaux            & 1 (mono) \\
Taille de trame   & 20 à 100\,ms par message (320 à 1600 échantillons) \\
Transport         & trames binaires WebSocket \\
\bottomrule
\end{tabular}
\end{center}

Les navigateurs captent en 44,1 ou 48\,kHz : la page doit donc sous-échantillonner à
16\,kHz. Utilisez un \texttt{AudioWorklet} plutôt que le \texttt{ScriptProcessorNode}
déprécié : il s'exécute hors du fil principal et ne produira pas de coupures lorsque
l'interface ERP est chargée. Rééchantillonner côté navigateur décharge également nos
serveurs d'inférence mutualisés, ce qui augmente directement le nombre de locuteurs
simultanés que nous pouvons servir.

Merci de ne pas envoyer de WebM, d'Opus ni de MP3. Le moteur consomme du PCM brut ; un
format conteneur devrait être décodé côté serveur et ajouterait de la latence.

\subsubsection*{Messages reçus}

Des trames texte, un objet JSON chacune.

\begin{lstlisting}
{ "type": "partial", "text": "bonjour je voudrais commander" }
\end{lstlisting}

\begin{lstlisting}
{ "type":  "final",
  "text":  "Bonjour, je voudrais commander trois cartons.",
  "start": 0.42,
  "end":   3.87 }
\end{lstlisting}

\begin{lstlisting}
{ "type": "error", "code": "ticket_expired", "message": "..." }
\end{lstlisting}

Affichez \texttt{partial} comme un texte provisoire — grisé ou en italique, à votre
convenance — et remplacez-le à l'arrivée du \texttt{final} correspondant. Plusieurs
partiels par seconde sont normaux, chacun remplaçant le précédent. Un \texttt{final} est
stable et ne sera pas révisé : seuls les finaux doivent être enregistrés. La socket se
ferme après un \texttt{error}.

\subsubsection*{Clore une session}

Envoyez une trame texte \texttt{\{"type": "close"\}} lorsque l'utilisateur arrête de
dicter, puis fermez la socket. Nous vidons l'audio encore en transit et émettons un
dernier \texttt{final} avant la fermeture. Couper la connexion fonctionne aussi, mais
risque de perdre la dernière phrase.

\section{Ce dont nous avons besoin de votre part}

\begin{askbox}
\begin{enumerate}
  \item \textbf{La ou les origines exactes du navigateur} qui ouvriront des sockets, par
        exemple \texttt{https://newapp.cynov.com}. Les WebSockets ne sont pas soumises à
        la politique de même origine : nous validons donc l'en-tête \texttt{Origin}
        contre une liste blanche lors de la poignée de main. Sans cela, un ticket dérobé
        fonctionnerait depuis n'importe quel site.
  \item \textbf{La concurrence attendue} : combien d'utilisateurs peuvent dicter
        simultanément en pointe. C'est notre véritable contrainte de capacité : le coût
        croît avec le nombre de locuteurs actifs, et les serveurs d'inférence sont
        uniquement CPU et partagés avec d'autres charges. Nous plafonnerons le nombre de
        sockets simultanées par clé API, et nous préférons fixer ce chiffre avec vous
        plutôt que le découvrir.
  \item \textbf{La destination de la transcription.} Si seul le navigateur en a besoin —
        l'utilisateur la lit et enregistre un formulaire — le flux ci-dessus est complet.
        Si votre \emph{backend} a besoin du texte, pour l'archiver dans une fiche,
        dites-le-nous maintenant : le faire remonter depuis le navigateur vers votre
        serveur est source de pertes et de conditions de course. Nous livrerions plutôt
        les finaux à un webhook que vous exposez, ou votre backend tiendrait la socket.
  \item \textbf{Un échantillon d'audio français représentatif}, si vous pouvez en
        partager un. Les chiffres de précision publiés pour ces modèles sont mesurés sur
        de la lecture. Le français spontané de terrain, avec interruptions, accents et
        bruit de fond, se comporte très différemment, et nous préférons mesurer sur vos
        données réelles plutôt que sur un banc d'essai.
\end{enumerate}
\end{askbox}

\section{Points ouverts de notre côté}

Listés par transparence. Aucun ne modifie ce que la section~\ref{sec:erp} vous demande.

\begin{itemize}
  \item \textbf{Latence.} Le modèle expose des tailles de bloc de 80, 160, 320, 560 et
        1120\,ms. Nous prévoyons 320 ou 560\,ms, ce qui maintient la latence de bout en
        bout sous la seconde. Des blocs plus courts paraissent plus immédiats et coûtent
        un peu de précision.
  \item \textbf{Débit sur CPU.} Notre parc est dépourvu de GPU. Nous mesurons le facteur
        temps réel sur le matériel réel avant de nous engager, car les bancs d'essai
        publiés du modèle sont uniquement sur GPU. C'est le principal risque sur le
        planning, et il conditionne la réponse au point~2 ci-dessus.
  \item \textbf{Comptage.} Le flux continu se facture en secondes d'audio et non en
        requêtes. Notre suivi de consommation par clé est étendu en conséquence.
\end{itemize}

\section{Synthèse du contrat}

\begin{center}
\renewcommand{\arraystretch}{1.45}
\begin{tabular}{@{}p{7.1cm}p{7.1cm}@{}}
\toprule
{\color{accent}\bfseries Équipe ERP} & {\color{accent}\bfseries Plateforme d'inférence} \\
\midrule
Conserver la clé API uniquement côté serveur &
Émettre des tickets éphémères, à usage unique, de portée réduite \\
Émettre un ticket par session de dictée &
Valider le ticket et l'\texttt{Origin} à la poignée de main \\
Capter le micro, sous-échantillonner en PCM16 16\,kHz mono &
Transcrire en continu, diffuser partiels et finaux \\
Afficher les partiels, n'enregistrer que les finaux &
Imputer toute la consommation à la clé API émettrice \\
Fournir origines, concurrence attendue, échantillon audio &
Plafonner les sockets par clé, publier le contrat définitif \\
\bottomrule
\end{tabular}
\end{center}

\vspace{1.1em}
\noindent
Si quoi que ce soit ici ne correspond pas à la façon dont l'ERP est construit, merci de
nous le signaler maintenant plutôt qu'une fois développé. La forme du flux
d'authentification est imposée par la contrainte de la section~2 ; presque tout le reste
est négociable.

\end{document}
